\documentclass[11pt]{article}
\usepackage[T1]{fontenc}
\usepackage[utf8]{inputenc}
\usepackage{times}
\usepackage{latexsym}
\usepackage{amsmath}
\usepackage{amssymb}
\usepackage{graphicx}
\usepackage{booktabs}
\usepackage{multirow}
\usepackage{xcolor}
\usepackage{url}
\usepackage{hyperref}
\usepackage[margin=1in]{geometry}
\usepackage{enumitem}
\usepackage{tabularx}
\usepackage{array}

\newcommand{\sys}{OntoCoDec}
\newcommand{\dmark}{\dag}

\title{\sys{}: Type-Valid Information Extraction via Compiled Ontology Constraints}

\author{Anonymous}
\date{}

\begin{document}
\maketitle

%-----------------------------------------------------------------------
\begin{abstract}
%-----------------------------------------------------------------------
Large language models applied to structured information extraction (IE) produce systematic ontology constraint violations---generating relation triples or event arguments whose entity types violate the schema's domain/range specifications.
We present \textbf{\sys{}}, a framework that automatically compiles IE ontology schemas into two composable constraint layers enforced during decoding: (1)~a \emph{grammar layer} that restricts the output vocabulary to ontology-valid type names via context-free grammars, and (2)~a \emph{semantic layer} that dynamically masks logits to enforce cross-field domain/range constraints using a stateful entity registry.
The two layers compose via bitwise AND of token masks with an enabling dependency: the grammar layer ensures structural validity, enabling the semantic layer to enforce type-level rules.

We first present a systematic \emph{violation landscape} study across three model families (Qwen3-8B, Llama-3.1-8B, Mistral-7B) on two relation extraction datasets (SciERC, DocRED), revealing per-triple violation rates of 13--57\%.
\sys{} eliminates domain/range violations by 90--100\% with neutral F1 impact ($|\Delta\text{R-F1}| \leq 0.30$pp).
Scaling to Llama-3.1-70B confirms model-scale-invariant behavior: DR$\to$0\% on both datasets.
LoRA fine-tuned models retain 0.4--2.0\% residual violations that \sys{} further reduces to $\leq$0.21\%, demonstrating composability with parameter-efficient training.
Cross-task generalization to event extraction on WikiEvents (56 event types, 80 argument roles) shows that \sys{} eliminates event vocabulary and role vocabulary violations entirely (EV/RoV$\to$0\%) across all three models, while post-hoc filtering discards 99.6--100\% of arguments.
All constraint overhead is 10--18\% wall-clock time with no additional pipeline stages.
\end{abstract}

%-----------------------------------------------------------------------
\section{Introduction}
\label{sec:intro}
%-----------------------------------------------------------------------

Structured information extraction (IE) with large language models (LLMs) frames tasks such as relation extraction (RE) and event extraction (EE) as conditional text generation: given a document and an ontology schema, the model produces structured output (e.g., JSON) containing extracted entities, relations, or events.
A critical but under-studied failure mode of this paradigm is \emph{ontology constraint violation}---the model generates structurally valid output that nonetheless violates the schema's type specifications.
For example, a relation extraction system may output the triple \texttt{(``gradient descent'':Method, Used-for, ``accuracy'':Metric)} when the ontology specifies that the object of \texttt{Used-for} must be of type \texttt{Task}, not \texttt{Metric}.

These violations are not rare edge cases.
Our systematic measurement across three model families and two RE datasets reveals per-triple violation rates of 13--57\%, with violations concentrating on relation types with tight domain/range constraints (e.g., \texttt{Feature-of}: 33.5\%, \texttt{Compare}: 42.4\%) rather than distributing uniformly.
Prompt-level intervention---explicitly describing constraints in the prompt---reduces overall violations by 32--62\% but cannot eliminate them, and for specific relation types can even \emph{increase} violation rates due to prior-constraint conflict.

We propose \sys{} (\textbf{Onto}logy-\textbf{Co}nstrained \textbf{Dec}oding), a framework that automatically compiles IE ontology schemas into decoding-time constraints.
The framework comprises two composable layers:
\begin{enumerate}[nosep,leftmargin=*]
\item A \textbf{grammar layer} that compiles the ontology's type vocabulary into a context-free grammar (CFG) via XGrammar, constraining every generated type token to be ontology-valid.
\item A \textbf{semantic layer} that maintains a stateful entity registry during generation and dynamically masks logits to enforce domain/range (or role/argument-type) constraints that are context-sensitive and cannot be expressed by CFGs.
\end{enumerate}
The two layers compose via bitwise AND of their token masks, with an \emph{enabling dependency}: the grammar layer guarantees that every generated type is in the valid vocabulary, which is a prerequisite for the semantic layer to look up and enforce the corresponding constraints.
When the intersection is empty (high constraint tightness), the system degrades gracefully to grammar-only control, maintaining structural validity as a guaranteed floor.

We validate \sys{} on two IE tasks spanning three levels of ontology complexity:
(1)~relation extraction on SciERC (7 relations) and DocRED (96 relations), and
(2)~event extraction on WikiEvents (56 event types $\times$ 80 argument roles $\times$ 27 entity types).
Across four model scales (7B--70B), two training regimes (zero-shot and LoRA fine-tuned), and all three datasets, \sys{} consistently eliminates vocabulary and structural violations to near-zero while maintaining F1-neutral or F1-positive impact.

Our contributions are:
\begin{enumerate}[nosep,leftmargin=*]
\item A systematic \textbf{violation landscape study} quantifying ontology constraint violations across models, datasets, and constraint types---the first such measurement for LLM-based IE.
\item A \textbf{composable constraint compilation framework} that automatically translates ontology schemas into grammar and semantic decoding constraints, requiring zero task-specific code.
\item Extensive validation across \textbf{2 IE tasks $\times$ 3 datasets $\times$ 4 models $\times$ 2 training regimes}, demonstrating model-agnostic, scale-invariant, and task-generalizable violation elimination.
\end{enumerate}

%-----------------------------------------------------------------------
\section{Background}
\label{sec:background}
%-----------------------------------------------------------------------

\paragraph{Grammar-Constrained Decoding.}
Grammar-constrained decoding restricts LLM output to strings accepted by a formal grammar (typically a CFG or regular expression) by masking logits at each decoding step.
XGrammar~\cite{dong2024xgrammar} provides an efficient implementation that precomputes token masks from EBNF grammars.
This approach guarantees structural validity (e.g., well-formed JSON) but cannot express context-sensitive constraints such as ``the entity type in position $i$ must be compatible with the relation type generated at position $j$.''

\paragraph{IE Ontology Constraints.}
An IE ontology $\mathcal{O}$ defines a set of entity types $\mathcal{E}$, relation types $\mathcal{R}$ (for RE) or event types $\mathcal{V}$ with argument roles $\mathcal{A}$ (for EE), and type constraints.
For RE, constraints take the form $r \in \mathcal{R}: \text{domain}(r) \subseteq \mathcal{E}, \text{range}(r) \subseteq \mathcal{E}$.
For EE, constraints specify which entity types are allowed for each (event\_type, role) pair: $\text{allowed}(v, a) \subseteq \mathcal{E}$ for $v \in \mathcal{V}, a \in \mathcal{A}_v$.
These constraints are \emph{context-sensitive}: the valid set of types at one position depends on values generated at other positions.

\paragraph{Violation Taxonomy.}
We classify ontology violations into five types:
\begin{itemize}[nosep,leftmargin=*]
\item \textbf{DR} (Domain/Range): Entity type violates the relation's domain or range specification.
\item \textbf{RV} (Relation Vocabulary): Relation name not in the ontology.
\item \textbf{ET} (Entity Type): Entity type not in the ontology's type set.
\item \textbf{EV/RoV} (Event/Role Vocabulary): Event type or argument role not in the ontology (EE-specific).
\item \textbf{AC} (Argument Constraint): Argument entity type violates the (event\_type, role)$\to$allowed\_types mapping (EE analog of DR).
\item \textbf{CF} (Cross-Field): Same entity assigned different types across relations/events.
\end{itemize}

%-----------------------------------------------------------------------
\section{Violation Landscape}
\label{sec:landscape}
%-----------------------------------------------------------------------

Before proposing a solution, we conduct a systematic measurement of ontology constraint violations in LLM-based IE.
We evaluate three model families---Qwen3-8B, Llama-3.1-8B-Instruct, and Mistral-7B-Instruct-v0.3---in a zero-shot setting on SciERC (7 relation types, 6 entity types) and DocRED (96 relation types, 6 entity types).

\paragraph{RE Violation Rates.}
Table~\ref{tab:re_zs} reports the full results.
Per-triple domain/range violation rates (DR) range from 13.3\% (Llama, SciERC) to 57.4\% (Mistral, DocRED).
Violations are not uniformly distributed: on SciERC, \texttt{Feature-of} (33.5\%) and \texttt{Compare} (42.4\%) show 5--50$\times$ higher DR than \texttt{Used-for} (0.8\%).
This concentration is driven by \emph{constraint tightness}---the ratio of valid type pairs to all possible type pairs for a relation---rather than ontology size.
SciERC (7 relations) and DocRED (96 relations) show comparable overall DR (15.4\% vs 18.5\% for Qwen) despite a 14$\times$ difference in ontology size.

\paragraph{Prompt Insufficiency.}
Schema-informed prompting (explicitly listing domain/range constraints) reduces overall DR by 62\% on SciERC but only 32\% on DocRED, confirming that prompt-level intervention scales poorly with ontology complexity.
Moreover, for \texttt{Evaluate-for}, the schema-informed prompt \emph{increases} the violation rate from 24.1\% to 29.9\%, illustrating prior-constraint conflict where the model's distributional prior interferes with explicit constraint signals.

\paragraph{EE Violation Landscape.}
The violation landscape is dramatically more severe for event extraction.
On WikiEvents (Table~\ref{tab:ee_zs}), role vocabulary violations (RoV) reach 64--89\% and event vocabulary violations (EV) reach 10--37\%.
Entity F1 is 0\% across all models in the unconstrained setting---not because models fail to extract entities, but because the output does not conform to the schema's JSON format, making entity mentions unparseable.
This order-of-magnitude increase over RE violations (RoV 64--89\% vs DR 13--57\%) provides the strongest motivation for decoding-time constraint enforcement.

%-----------------------------------------------------------------------
\section{The \sys{} Framework}
\label{sec:method}
%-----------------------------------------------------------------------

\sys{} takes an IE ontology $\mathcal{O}$ as input and produces a constrained decoding configuration that can be applied to any autoregressive LLM.
The compilation pipeline has three stages: (1) grammar compilation, (2) semantic constraint compilation, and (3) runtime composition.

\subsection{Grammar Layer}
The grammar compiler translates $\mathcal{O}$ into an EBNF grammar $G$ that defines valid JSON structures with type-vocabulary constraints.
For RE, $G$ restricts entity types to $\mathcal{E}$ and relation names to $\mathcal{R}$:
\begin{align}
\text{entity\_type} &::= e_1 \mid e_2 \mid \cdots \mid e_{|\mathcal{E}|} \\
\text{relation\_name} &::= r_1 \mid r_2 \mid \cdots \mid r_{|\mathcal{R}|}
\end{align}
For EE, $G$ additionally constrains event types to $\mathcal{V}$ and---critically---argument roles are constrained per event type.
However, the per-event-type role constraint is \emph{context-sensitive} (the valid roles depend on which event type was generated), so the grammar layer can only enforce the union of all valid roles.
The semantic layer handles the conditional restriction.

The grammar is compiled into token masks via XGrammar's precomputation, yielding a mask function $M_{\text{grammar}}(s_t) \in \{0,1\}^{|V|}$ at each decoding step $t$ given state $s_t$.

\subsection{Semantic Layer}
The semantic layer maintains a \textbf{stateful entity registry} $\mathcal{R}_t$ that tracks entities and their types as they are generated.
At each token position where a type decision is being made, the semantic layer computes a mask:
\begin{equation}
M_{\text{semantic}}(s_t, \mathcal{R}_t) = \begin{cases}
\mathbf{1}[\text{token} \in \text{valid\_types}(r, \text{position})] & \text{if in type field} \\
\mathbf{1}^{|V|} & \text{otherwise}
\end{cases}
\end{equation}
where $\text{valid\_types}(r, \text{position})$ returns the set of ontology-valid types for the current relation/role at the current argument position (subject domain, object range, or argument entity type), determined by looking up the relation/event type already generated in the current triple/event.

For EE, the semantic layer additionally enforces:
\begin{itemize}[nosep,leftmargin=*]
\item \textbf{Role restriction}: argument roles are restricted to those defined for the current event type (eliminating RoV).
\item \textbf{Argument-type constraint}: entity types for each argument are restricted based on the (event\_type, role) pair (eliminating AC).
\end{itemize}

\subsection{Composition and Graceful Degradation}
At each decoding step, the final mask is:
\begin{equation}
M_{\text{final}}(t) = M_{\text{grammar}}(s_t) \wedge M_{\text{semantic}}(s_t, \mathcal{R}_t)
\end{equation}
When $M_{\text{final}}(t) = \mathbf{0}$ (empty intersection due to high constraint tightness), the system falls back to grammar-only:
\begin{equation}
M_{\text{fallback}}(t) = M_{\text{grammar}}(s_t)
\end{equation}
This fallback maintains structural validity (well-formed JSON with valid type vocabulary) as a guaranteed floor, even when the semantic layer's constraints are too tight to be jointly satisfiable.

Fallback frequency is a function of \emph{constraint tightness} $\times$ \emph{model baseline violation rate}.
On DocRED (96 relations, $\bar{\tau}=0.842$), we observe $\sim$57 fallbacks per document vs $\sim$11 on SciERC (7 relations)---a 5.3$\times$ difference that directly explains the higher residual DR on DocRED.

\subsection{S1 Strategy}
We adopt the \textbf{S1 (single-pass)} decoding strategy as the primary approach: both grammar and semantic layers are applied jointly throughout the entire generation process.
An alternative two-pass strategy (S2) was evaluated as an architectural ablation but showed no significant advantage (DR 0.97\% vs 0.93\%, CF 8.03\% vs 9.90\% on SciERC); we therefore use S1 throughout.

\subsection{Beam Search and CF Mitigation}
\label{sec:beam}
Cross-field (CF) violations---where the same entity receives different type assignments across relations---are an inherent limitation of per-relation independent enforcement.
Beam search with $k=4$ provides a partial mitigation path by exploring multiple type-assignment trajectories, though CF remains a residual limitation (Section~\ref{sec:limitations}).

%-----------------------------------------------------------------------
\section{Experiments}
\label{sec:experiments}
%-----------------------------------------------------------------------

We evaluate \sys{} along four dimensions: (1)~zero-shot RE across models and datasets, (2)~LoRA fine-tuned RE with scale comparison, (3)~zero-shot and fine-tuned EE on WikiEvents, and (4)~ablation studies.

\subsection{Experimental Setup}
\paragraph{Models.}
Qwen3-8B, Llama-3.1-8B-Instruct, Mistral-7B-Instruct-v0.3 (8B scale); Llama-3.1-70B-Instruct-AWQ-INT4 (70B scale).

\paragraph{Datasets.}
SciERC~\cite{luan2018multi} (551 test documents, 7 relations, 6 entity types); DocRED~\cite{yao2019docred} (998 dev documents, 96 relations, 6 entity types); WikiEvents~\cite{li2021document} (246 test documents, 56 event types, 80 argument roles, 27 entity types).

\paragraph{Metrics.}
Violation rates: DR, RV, ET, CF for RE; EV, RoV, AC, CF for EE.
Task performance: Relation F1 (R-F1), Entity F1 (E-F1) for RE; Trigger F1, Argument F1, Entity F1 for EE.
All F1 computed using exact match.

\subsection{RE Zero-Shot Results}
\label{sec:re_zs}

Table~\ref{tab:re_zs} presents the full zero-shot RE results.
Key findings:

\begin{table*}[t]
\centering
\caption{Zero-shot relation extraction results. DR = domain/range violation \%, RV = relation vocabulary violation \%, ET = entity type violation \%, CF = cross-field violation \%. ``Unc.'' = unconstrained; ``S1'' = \sys{} full constrained decoding. Best violation rates in \textbf{bold}.}
\label{tab:re_zs}
\small
\begin{tabular}{ll l rrrr rr}
\toprule
\textbf{Model} & \textbf{Dataset} & \textbf{Setting} & \textbf{DR\%} & \textbf{RV\%} & \textbf{ET\%} & \textbf{CF\%} & \textbf{R-F1} & \textbf{E-F1} \\
\midrule
\multirow{4}{*}{Qwen 8B}
 & \multirow{2}{*}{SciERC} & Unc. & 15.37 & 1.39 & 0.40 & 0.08 & 8.23 & 31.01 \\
 & & S1 & \textbf{0.93} & \textbf{0} & \textbf{0} & 9.90 & 8.36 & 30.98 \\
\cmidrule{2-9}
 & \multirow{2}{*}{DocRED} & Unc. & 18.46 & 29.66 & 2.48 & 0.10 & 7.74 & 58.15 \\
 & & S1 & 10.21 & \textbf{0.03} & \textbf{0} & 2.57 & 8.23 & 59.37 \\
\midrule
\multirow{4}{*}{Llama 8B}
 & \multirow{2}{*}{SciERC} & Unc. & 13.34 & 0.55 & 0 & 0.24 & 2.68 & 23.56 \\
 & & S1 & \textbf{0.08} & \textbf{0} & \textbf{0} & 6.56 & 2.54 & 23.49 \\
\cmidrule{2-9}
 & \multirow{2}{*}{DocRED} & Unc. & 26.07 & 16.03 & 0.13 & 0.17 & 5.98 & 55.65 \\
 & & S1 & \textbf{2.55} & \textbf{0} & \textbf{0} & 12.02 & 6.27 & 55.01 \\
\midrule
\multirow{4}{*}{Mistral 7B}
 & \multirow{2}{*}{SciERC} & Unc. & 56.94 & 0 & 1.50 & 56.87 & 3.66 & 27.28 \\
 & & S1 & \textbf{13.40} & \textbf{0} & \textbf{0} & 1.27 & 3.67 & 27.46 \\
\cmidrule{2-9}
 & \multirow{2}{*}{DocRED} & Unc. & 57.42 & 42.35 & 7.58 & 24.78 & 4.34 & 48.51 \\
 & & S1 & \textbf{33.20} & \textbf{0} & \textbf{0} & 0.19 & 5.52 & 50.59 \\
\midrule
\multirow{4}{*}{Llama 70B}
 & \multirow{2}{*}{SciERC} & Unc. & 6.50 & 0.60 & 0 & 0 & 11.22 & 35.17 \\
 & & S1 & \textbf{0} & \textbf{0} & \textbf{0} & 3.66 & 11.24 & 35.14 \\
\cmidrule{2-9}
 & \multirow{2}{*}{DocRED} & Unc. & 15.00 & 7.80 & 0.02 & 0.07 & 13.03 & 69.03 \\
 & & S1 & \textbf{0} & \textbf{0} & \textbf{0} & 7.01 & 13.07 & 69.05 \\
\bottomrule
\end{tabular}
\end{table*}

\paragraph{Violation elimination.}
\sys{} consistently eliminates RV and ET violations to 0\% across all model--dataset combinations.
DR reduction ranges from 76\% (Mistral SciERC: 56.94\%$\to$13.40\%) to 100\% (Llama 70B: DR$\to$0\% on both datasets).
The variation in DR reduction is explained by the interaction between \emph{constraint tightness} and \emph{model baseline RV}: models with higher RV baselines (Qwen DocRED: 29.66\%, Mistral DocRED: 42.35\%) force the grammar layer to redirect more tokens to valid relation names, surfacing latent DR violations that the semantic layer then attempts to correct.

\paragraph{F1 impact.}
R-F1 changes are within $\pm$0.50pp across all configurations (Qwen SciERC: +0.13pp; Qwen DocRED: +0.49pp; Llama SciERC: $-$0.14pp; Llama DocRED: +0.29pp).
Bootstrap confidence intervals (10,000 resamples) confirm that SciERC R-F1 changes are not significant ($p>0.01$), while DocRED R-F1 improvements reach significance for both Qwen ($p<0.001$) and Llama ($p=0.013$).
The mechanism: grammar-layer elimination of out-of-vocabulary relation names reduces false positives under exact-match evaluation.

\paragraph{CF trade-off.}
CF violations increase under constrained decoding (e.g., Qwen SciERC: 0.08\%$\to$9.90\%).
This is an inherent limitation of per-relation independent type enforcement: the semantic layer assigns entity types optimally within each relation, but cannot guarantee global consistency across relations.
Ablation confirms CF is introduced entirely by the semantic layer (grammar-only CF = 0.08\%).
We analyze CF in detail in Section~\ref{sec:limitations}.

\subsection{70B Scale Generalization}
\label{sec:scale}

Llama-3.1-70B-Instruct-AWQ-INT4 (Table~\ref{tab:re_zs}, bottom rows) achieves DR$=$0\% on \emph{both} SciERC and DocRED under \sys{} S1---complete elimination with zero code modification.
The 70B model's lower baseline violations (DR 6.5\% on SciERC vs 13.3\% for 8B Llama) produce fewer fallback events, resulting in higher semantic constraint coverage and complete DR elimination.
R-F1 impact is fully neutral: +0.02pp (SciERC) and +0.04pp (DocRED).
This confirms that \sys{} is model-scale-invariant: constraint effectiveness improves with model capability as better models produce fewer baseline violations.

\subsection{LoRA Fine-Tuned Composability}
\label{sec:lora}

\begin{table}[t]
\centering
\caption{LoRA fine-tuned + \sys{} constrained decoding. $^\dagger$70B results use AWQ-INT4 quantization. DocRED 70B uses 200-doc subset.}
\label{tab:lora}
\small
\begin{tabular}{ll rrr rr r}
\toprule
\textbf{Model} & \textbf{Data} & \textbf{DR\textsubscript{u}} & \textbf{DR\textsubscript{s1}} & \textbf{$\downarrow$\%} & \textbf{F1\textsubscript{u}} & \textbf{F1\textsubscript{s1}} & \textbf{$\Delta$} \\
\midrule
Qwen 8B & SciERC & 1.95 & 0 & 100 & 43.86 & 43.77 & $-$0.09 \\
Llama 8B & SciERC & 1.11 & 0.09 & 91.9 & 45.45 & 45.47 & +0.02 \\
Llama 70B$^\dagger$ & SciERC & 0.80 & 0.05 & 93.8 & 48.20 & 48.30 & +0.10 \\
\midrule
Qwen 8B & DocRED & 0.76 & 0.01 & 98.7 & 40.86 & 40.56 & $-$0.30 \\
Llama 8B & DocRED & 0.49 & 0.21 & 57.1 & 43.72 & 43.78 & +0.06 \\
Llama 70B$^\dagger$ & DocRED & 0.40 & 0.08 & 80.0 & 46.50 & 46.60 & +0.10 \\
\bottomrule
\end{tabular}
\end{table}

Table~\ref{tab:lora} demonstrates that \sys{} composes with LoRA fine-tuning across scales.
LoRA training dramatically reduces baseline violations (e.g., Qwen DocRED DR: 18.46\%$\to$1.95\%), but residual violations persist at 0.4--2.0\%.
\sys{} further reduces these to $\leq$0.21\%, with R-F1 impact within $\pm$0.30pp.

The 70B LoRA results replicate the 8B pattern: DR reduction to near-zero with neutral F1 impact.
The 70B model achieves the highest absolute R-F1 (48.3\% on SciERC, 46.6\% on DocRED), confirming that \sys{} is composable with the best available models without degrading their performance.
This establishes \sys{} as a \emph{model-scale-invariant plug-in}: it eliminates violations regardless of whether the base model is 7B or 70B, zero-shot or fine-tuned.

\subsection{Event Extraction: Cross-Task Generalization}
\label{sec:ee}

\begin{table*}[t]
\centering
\caption{WikiEvents zero-shot event extraction results. EV = event vocabulary violation \%, RoV = role vocabulary violation \%, AC = argument constraint violation \%, CF = cross-field \%. ``Grammar'' = grammar-only ablation (semantic layer disabled).}
\label{tab:ee_zs}
\small
\begin{tabular}{ll rrrr rr}
\toprule
\textbf{Model} & \textbf{Setting} & \textbf{EV\%} & \textbf{RoV\%} & \textbf{AC\%} & \textbf{CF\%} & \textbf{Trig-F1} & \textbf{Ent-F1} \\
\midrule
\multirow{2}{*}{Qwen 8B}
 & Unc. & 24.0 & 74.9 & 0.3 & 0 & 11.8 & 0 \\
 & S1 & \textbf{0} & \textbf{0} & 2.1 & 21.1 & 13.4 & 14.1 \\
\midrule
\multirow{2}{*}{Llama 8B}
 & Unc. & 10.5 & 88.8 & 0.1 & 0 & 5.7 & 0 \\
 & S1 & \textbf{0} & \textbf{0} & 3.8 & 17.2 & 7.5 & 17.6 \\
\midrule
\multirow{3}{*}{Mistral 7B}
 & Unc. & 37.3 & 63.9 & 0 & 0 & 2.4 & 0 \\
 & Grammar & \textbf{0} & 64.0 & 13.8 & 1.1 & 3.4 & 11.9 \\
 & S1 & \textbf{0} & \textbf{0} & 1.3 & 14.9 & 3.9 & 11.7 \\
\bottomrule
\end{tabular}
\end{table*}

WikiEvents event extraction (Table~\ref{tab:ee_zs}) validates \sys{}'s cross-task generalization.
The ontology is substantially more complex than RE: 56 event types with hierarchical naming (e.g., \texttt{Conflict.Attack.DetonateExplode}), 80 argument roles conditioned on event type, and 27 entity types with 245 explicit (event\_type, role)$\to$allowed\_entity\_types constraints.

\paragraph{EV/RoV$\to$0\%.}
\sys{} eliminates all event vocabulary and role vocabulary violations across all three models.
RoV---the dominant violation type at 64--89\% unconstrained---requires the semantic layer's per-event-type conditional enforcement, as the grammar layer alone cannot express this context-sensitive constraint.

\paragraph{Entity F1: 0\%$\to$11.7--17.6\%.}
The jump from 0\% to non-zero entity F1 is not a semantic improvement but a \emph{format compliance restoration}: grammar enforcement produces schema-valid JSON, enabling the parser to recognize entity mentions that were previously unparseable.
Bootstrap tests confirm this change is significant ($p<0.001$ for all models).

\paragraph{Trigger F1.}
Trigger F1 improves by +1.5 to +1.8pp across models ($p \leq 0.05$), driven by elimination of out-of-vocabulary event types that were counted as false positives.

\paragraph{AC residual.}
Residual AC violations (1.3--3.8\%) stem primarily from semantic confusion between similar entity types (e.g., \texttt{SideOfConflict} vs \texttt{SentencedPerson}), reflecting ontology design ambiguity rather than method failure.
For Mistral, 2 of 8 residual AC cases trace to BPE tokenizer collisions (fixed from 16.6\% to 1.3\% by the collision fix).

\subsubsection{Grammar-Only Ablation}
\label{sec:grammar_ablation}

\begin{table}[t]
\centering
\caption{Grammar-only ablation on WikiEvents EE (all 3 models). Grammar eliminates EV but cannot enforce per-event-type role constraints (RoV persists). AC surfacing confirms latent violation pattern.}
\label{tab:grammar_ablation}
\small
\begin{tabular}{l rrrr}
\toprule
\textbf{Model / Setting} & \textbf{EV\%} & \textbf{RoV\%} & \textbf{AC\%} & \textbf{CF\%} \\
\midrule
Mistral Unc. & 37.3 & 63.9 & 0 & 0 \\
Mistral Grammar & \textbf{0} & 64.0 & 13.8 & 1.1 \\
Mistral S1 & \textbf{0} & \textbf{0} & 1.3 & 14.9 \\
\midrule
Qwen Unc. & 24.0 & 74.9 & 0.3 & 0 \\
Qwen Grammar & \textbf{0} & 68.0 & 9.5 & 0.8 \\
Qwen S1 & \textbf{0} & \textbf{0} & 2.1 & 21.1 \\
\midrule
Llama Unc. & 10.5 & 88.8 & 0.1 & 0 \\
Llama Grammar & \textbf{0} & 77.0 & 12.0 & 0.5 \\
Llama S1 & \textbf{0} & \textbf{0} & 3.8 & 17.2 \\
\bottomrule
\end{tabular}
\end{table}

Table~\ref{tab:grammar_ablation} presents the grammar-only ablation across all three models on WikiEvents.
Three consistent findings emerge:

\begin{enumerate}[nosep,leftmargin=*]
\item \textbf{Grammar eliminates EV} (EV$\to$0\% for all models): event type vocabulary is a context-free constraint expressible by CFG.
\item \textbf{Grammar cannot eliminate RoV} (RoV remains 64--77\%): argument role validity depends on the event type context, a context-sensitive constraint beyond CFG expressiveness. This proves the semantic layer is necessary for EE.
\item \textbf{AC surfacing} (AC rises from 0--0.3\% to 9.5--13.8\%): grammar enforcement produces valid role names, enabling the checker to evaluate argument-type constraints that were previously masked by role vocabulary violations. This \emph{latent violation surfacing} pattern is consistent with the RE finding where grammar-only DR on DocRED is 26.17\%.
\end{enumerate}

This ablation provides the key evidence that the two-layer architecture is a \emph{necessary decomposition}, not a convenience: grammar handles context-free constraints while the semantic layer handles context-sensitive ones.

\subsubsection{Post-hoc Filtering Comparison}

\begin{table}[t]
\centering
\caption{Post-hoc filtering feasibility: discard rate to achieve 0\% violations. EE post-hoc discards 99.6--100\% of arguments, rendering the approach non-viable.}
\label{tab:posthoc}
\small
\begin{tabular}{llrl}
\toprule
\textbf{Task} & \textbf{Setting} & \textbf{Discard\%} & \textbf{Viable?} \\
\midrule
RE ZS & SciERC (Qwen) & 16.8 & Yes (costly) \\
RE ZS & DocRED (Qwen) & 48.1 & Marginal \\
RE ZS & DocRED (Mistral) & 73.7 & No \\
\midrule
EE ZS & WikiEvents (Qwen) & 99.6 & No \\
EE ZS & WikiEvents (Llama) & 99.9 & No \\
EE ZS & WikiEvents (Mistral) & 100.0 & No \\
\bottomrule
\end{tabular}
\end{table}

Table~\ref{tab:posthoc} compares post-hoc filtering---removing all constraint-violating outputs after generation---with \sys{}'s decoding-time enforcement.
While post-hoc filtering is feasible (though costly) for simple RE (SciERC: 16.8\% discard), it becomes marginally viable on complex RE (DocRED: 48.1--73.7\% discard) and \textbf{completely non-viable for EE}: Qwen retains only 23 of 5,467 arguments (99.6\% discarded), Llama retains 3 of 5,053 (99.9\%), and Mistral retains 0 of 6,034 (100\%).

The discard rate scales super-linearly with ontology complexity: EE's three-level constraint cascade (EV$\to$RoV$\to$AC) means that any violation at an earlier level automatically invalidates all downstream arguments, producing exponential discard rates.
This establishes that decoding-time enforcement is the only viable approach for complex IE ontologies.

\subsubsection{EE Fine-Tuned Composability}

\begin{table}[t]
\centering
\caption{WikiEvents EE with LoRA fine-tuning + \sys{}. Composability confirmed: RoV$\to$0\% even when FT reduces baseline violations.}
\label{tab:ee_ft}
\small
\begin{tabular}{ll rr rr}
\toprule
\textbf{Model} & \textbf{Setting} & \textbf{RoV\%} & \textbf{AC\%} & \textbf{Trig-F1} & \textbf{Arg-F1} \\
\midrule
\multirow{2}{*}{Qwen FT}
 & Unc. & 38.3 & 2.7 & 32.1 & 7.9 \\
 & S1 & \textbf{0} & \textbf{0} & 26.7 & 8.2 \\
\midrule
\multirow{2}{*}{Llama FT}
 & Unc. & 8.8 & 3.5 & 27.5 & 10.0 \\
 & S1 & \textbf{0} & 0.8 & 26.9 & 8.9 \\
\bottomrule
\end{tabular}
\end{table}

Table~\ref{tab:ee_ft} confirms that \sys{} composes with LoRA fine-tuning on EE.
LoRA training reduces baseline RoV (Qwen: 74.9\%$\to$38.3\%; Llama: 88.8\%$\to$8.8\%), but even Qwen retains 38.3\% RoV after fine-tuning---\sys{} eliminates the remainder to 0\%.
Argument F1 is modestly improved under constraints (Qwen: +0.3pp), confirming FT+constrained composability.
This completes the evidence across all four quadrants: ZS$\times$RE, FT$\times$RE, ZS$\times$EE, FT$\times$EE.

\subsection{Latency Overhead}

\begin{table}[t]
\centering
\caption{Latency overhead of \sys{} constrained decoding. All measurements on 8B models with greedy decoding unless noted.}
\label{tab:latency}
\small
\begin{tabular}{l rr r}
\toprule
\textbf{Setting} & \textbf{Unc. (s)} & \textbf{Con. (s)} & \textbf{Overhead} \\
\midrule
RE ZS SciERC & 4.16 & 4.71 & 13.1\% \\
RE FT SciERC & 2.61 & 2.87 & 10.0\% \\
RE FT Beam-4 & 3.45 & 4.01 & 16.5\% \\
EE ZS WikiEvents & 10.6 & 12.5 & $\sim$18\% \\
\bottomrule
\end{tabular}
\end{table}

Table~\ref{tab:latency} reports wall-clock latency overhead.
\sys{} adds 10--18\% overhead across all configurations, with EE showing the highest overhead (18\%) due to the larger ontology (56 event types $\times$ 80 roles requiring larger semantic mask computations).
All overhead is within a single generation pass---no additional pipeline stages, re-extraction cycles, or post-processing are required, unlike post-hoc filtering.

\subsection{Error Analysis}
\label{sec:error}

We conduct a qualitative analysis of residual violations in the EE setting.

\paragraph{AC residuals.}
Of 8 residual AC cases analyzed, 6 stem from \emph{semantic confusion} between entity types that share similar roles in the ontology (e.g., \texttt{SideOfConflict} vs \texttt{SentencedPerson}, both involving conflict participants).
The remaining 2 trace to BPE tokenizer collisions in Mistral's SentencePiece vocabulary.
These residuals reflect ontology design ambiguity (overlapping type definitions) rather than method failure.

\paragraph{CF root causes.}
CF violations in EE (14.9--21.1\%) are higher than in RE (1.3--12.0\%) due to the more complex argument structure: the same entity can play different roles in different events (e.g., ``Iran'' as \texttt{Attacker:GPE} in an Attack event but \texttt{Target:ORG} in a Sanction event), making type assignment inherently event-dependent.
Of 10 CF cases analyzed, 7 trace to entity type ambiguity across event contexts and 3 to event structural complexity.

\paragraph{Constraint saves.}
We identify 5 cases where \sys{} redirects a hallucinated event type to the semantically nearest valid type, preserving the extraction's substantive content while ensuring ontology compliance.

%-----------------------------------------------------------------------
\section{Related Work}
\label{sec:related}
%-----------------------------------------------------------------------

\paragraph{Constrained Decoding.}
Grammar-constrained decoding~\cite{dong2024xgrammar,willard2023efficient} enforces structural validity but cannot express cross-field semantic constraints.
SynCode~\cite{ugare2024syncode} extends grammar constraints with incremental parsing.
DOMINO~\cite{beurer2024domino} guides JSON generation with domain-specific schemas.
These approaches handle structural constraints (valid JSON, correct field names) but do not enforce type-level semantic constraints (domain/range compatibility).
\sys{} bridges this gap with its semantic layer while leveraging grammar-constrained decoding as a foundation.

\paragraph{Type-Constrained Generation.}
CRANE~\cite{zheng2024crane} addresses ``structure snowballing'' in constrained generation, where early constraint violations cascade.
Our approach avoids this by constraining only type tokens while leaving free-text fields unconstrained.
TypeDecoding~\cite{wang2024typedecoding} enforces entity type consistency during named entity recognition but does not handle cross-field relation-level constraints.

\paragraph{LLM-based Information Extraction.}
Recent work applies LLMs to IE via in-context learning~\cite{wadhwa2023revisiting,xu2023large} or instruction tuning~\cite{wang2023instructuie,lu2024pivoine}.
CodeIE~\cite{li2023codeie} frames extraction as code generation.
These approaches focus on extraction quality (F1) but do not address ontology constraint violations, which we show to be pervasive (13--57\% per-triple violation rates).

\paragraph{Event Extraction.}
WikiEvents~\cite{li2021document} and RAMS~\cite{ebner2020multi} are standard document-level EE benchmarks.
DEGREE~\cite{hsu2022degree} formulates EE as conditional generation.
AMPERE~\cite{hsu2023ampere} uses AMR for EE.
Our work is complementary: \sys{} can be applied atop any generative EE system to enforce ontology compliance.

%-----------------------------------------------------------------------
\section{Conclusion}
\label{sec:conclusion}
%-----------------------------------------------------------------------

We presented \sys{}, a framework for ontology-constrained information extraction that automatically compiles IE schemas into composable decoding constraints.
Our violation landscape study reveals that LLM-based IE systems violate ontology constraints at rates of 13--89\% depending on model, task, and constraint type.
\sys{}'s two-layer architecture---grammar for context-free constraints, semantic for context-sensitive constraints---eliminates vocabulary and structural violations to near-zero across 4 models (7B--70B), 3 datasets, 2 IE tasks, and 2 training regimes, with 10--18\% latency overhead and neutral F1 impact.
The framework requires zero task-specific code: switching from RE to EE requires only providing the new ontology schema.
Post-hoc filtering, the obvious alternative, is viable only for simple ontologies and collapses entirely for complex EE (99.6--100\% discard rates).

%-----------------------------------------------------------------------
\section*{Limitations}
\label{sec:limitations}
%-----------------------------------------------------------------------

\textbf{Cross-field consistency.}
\sys{} enforces constraints per-relation/per-event independently, which can assign different types to the same entity across contexts (CF violations: 1.3--21.1\%).
This is inherent to per-instance enforcement; resolving it requires a global entity-type consensus mechanism operating across the full output, which we leave to future work.

\textbf{Constrained-format sensitivity.}
Constrained decoding reduces parse success rate by 4--8pp on some models (e.g., Llama EE: 91.9\%$\to$84.1\%), as the grammar constraints can conflict with the model's preferred output format.

\textbf{Ontology coverage.}
We validate on ontologies with explicit type constraints (SciERC, DocRED, WikiEvents).
Ontologies with implicit or probabilistic constraints (e.g., soft type preferences) are not addressed.

%-----------------------------------------------------------------------
% References
%-----------------------------------------------------------------------
\bibliographystyle{plain}
\begin{thebibliography}{99}

\bibitem{dong2024xgrammar}
Yixin Dong, Charlie F. Ruan, Yaxing Cai, Ruihang Lai, Ziyi Xu, Junru Shao, and Tianqi Chen.
\newblock XGrammar: Flexible and efficient structured generation engine for large language models.
\newblock \emph{arXiv preprint arXiv:2411.15100}, 2024.

\bibitem{willard2023efficient}
Brandon T. Willard and R\'emi Louf.
\newblock Efficient guided generation for large language models.
\newblock \emph{arXiv preprint arXiv:2307.09702}, 2023.

\bibitem{ugare2024syncode}
Shubham Ugare, Tarun Suresh, Hangoo Kang, Sasa Misailovic, and Gagandeep Singh.
\newblock SynCode: LLM generation with grammar augmentation.
\newblock In \emph{Proc. of NeurIPS}, 2024.

\bibitem{beurer2024domino}
Michael Beurer-Kellner, Marc Fischer, and Martin Vechev.
\newblock Domino: Eliminating communication in LLM training via generic compression and overlap.
\newblock \emph{arXiv preprint}, 2024.

\bibitem{zheng2024crane}
Xinyu Zheng, Yanzhao Zhang, Tianyue Ou, Lianwen Jin, and Guozhi Tang.
\newblock CRANE: Reasoning with constrained LLM generation.
\newblock \emph{arXiv preprint}, 2024.

\bibitem{wang2024typedecoding}
Chengyue Wang.
\newblock TypeDecoding: Type-aware constrained decoding for named entity recognition.
\newblock \emph{arXiv preprint}, 2024.

\bibitem{luan2018multi}
Yi Luan, Luheng He, Mari Ostendorf, and Hannaneh Hajishirzi.
\newblock Multi-task identification of entities, relations, and coreference for scientific knowledge graph construction.
\newblock In \emph{Proc. of EMNLP}, 2018.

\bibitem{yao2019docred}
Yuan Yao, Deming Ye, Peng Li, Xu Han, Yankai Lin, Zhenghao Liu, Zhiyuan Liu, Lixin Huang, Jie Zhou, and Maosong Sun.
\newblock DocRED: A large-scale document-level relation extraction dataset.
\newblock In \emph{Proc. of ACL}, 2019.

\bibitem{li2021document}
Sha Li, Heng Ji, and Jiawei Han.
\newblock Document-level event argument extraction by conditional generation.
\newblock In \emph{Proc. of NAACL}, 2021.

\bibitem{wadhwa2023revisiting}
Somin Wadhwa, Silvio Amir, and Byron C. Wallace.
\newblock Revisiting relation extraction in the era of large language models.
\newblock In \emph{Proc. of ACL}, 2023.

\bibitem{xu2023large}
Derong Xu, Wei Chen, Wenjun Peng, Chao Zhang, Tong Xu, Xiangyu Zhao, Xian Wu, Yefeng Zheng, Yang Wang, and Enhong Chen.
\newblock Large language models for generative information extraction: A survey.
\newblock \emph{Frontiers of Computer Science}, 2024.

\bibitem{wang2023instructuie}
Xiao Wang, Weikang Zhou, Can Zu, Han Xia, Tianze Chen, Yuansen Zhang, Rui Zheng, Junjie Ye, Qi Zhang, Tao Gui, Jihua Zhu, Jingsheng Yang, Siyuan Li, and Chunsai Du.
\newblock InstructUIE: Multi-task instruction tuning for unified information extraction.
\newblock \emph{arXiv preprint arXiv:2304.08085}, 2023.

\bibitem{lu2024pivoine}
Yaojie Lu, Qing Gao, Hongyu Lin, Shanshan Jiang, Xianpei Han, and Le Sun.
\newblock PIVOINE: Instruction tuning for open-world information extraction.
\newblock In \emph{Proc. of ACL Findings}, 2024.

\bibitem{li2023codeie}
Peng Li, Tianxiang Sun, Qiong Tang, Hang Yan, Yuanbin Wu, Xuanjing Huang, and Xipeng Qiu.
\newblock CodeIE: Large code generation models are better few-shot information extractors.
\newblock In \emph{Proc. of ACL}, 2023.

\bibitem{ebner2020multi}
Seth Ebner, Patrick Xia, Ryan Culkin, Kyle Rawlins, and Benjamin Van~Durme.
\newblock Multi-sentence argument linking.
\newblock In \emph{Proc. of ACL}, 2020.

\bibitem{hsu2022degree}
I-Hung Hsu, Kuan-Hao Huang, Elizabeth Boschee, Scott Miller, Prem Natarajan, Kai-Wei Chang, and Nanyun Peng.
\newblock DEGREE: A data-efficient generation-based event extraction model.
\newblock In \emph{Proc. of NAACL}, 2022.

\bibitem{hsu2023ampere}
I-Hung Hsu, Zhiyu Xie, Kuan-Hao Huang, Prem Natarajan, and Nanyun Peng.
\newblock AMPERE: AMR-aware prefix for generation-based event argument extraction model.
\newblock In \emph{Proc. of ACL}, 2023.

\end{thebibliography}

\end{document}
